\textbf{Phần 1. Trắc nghiệm lựa chọn}
\setcounterpageref{ex}{0}

\begin{ex} %[2H3V18]
Một can dầu chứa $18\text{ L}$. Nếu rót thêm vào can $5\text{ L}$ nữa thì can đầy. Hỏi dung tích của chiếc can đó là bao nhiêu?
\choice
{$21\text{ L}$}
{\True $23\text{ L}$}
{$22\text{ L}$}
{$24\text{ L}$}
\loigiai{
Dung tích can: $18 + 5 = 23\text{ (L)}$.
\textbf{Đáp án của câu hỏi là B}
}
\end{ex}

\begin{ex} %[2H3V18]
Túi thứ nhất nặng $15\text{ kg}$, túi thứ hai nặng $18\text{ kg}$. Người ta lấy bớt ở túi thứ hai $4\text{ kg}$. Hỏi lúc này cả hai túi nặng bao nhiêu ki-lô-gam?
\choice
{\True $29\text{ kg}$}
{$31\text{ kg}$}
{$28\text{ kg}$}
{$30\text{ kg}$}
\loigiai{
Túi thứ hai còn: $18 - 4 = 14\text{ (kg)}$. Cả hai túi nặng: $15 + 14 = 29\text{ (kg)}$.
\textbf{Đáp án của câu hỏi là A}
}
\end{ex}

\begin{ex} %[2H3V18]
Điền số thích hợp: $42\text{ kg} - 17\text{ kg} = 15\text{ kg} + \dots\text{ kg}$.
\choice
{$9$}
{$11$}
{\True $10$}
{$8$}
\loigiai{
Vế trái: $42 - 17 = 25\text{ kg}$. Vế phải: $25 - 15 = 10$. Vậy điền 10.
\textbf{Đáp án của câu hỏi là C}
}
\end{ex}

\begin{ex} %[2H3V18]
Có hai thùng nước: thùng $A$ chứa $35\text{ L}$, thùng $B$ chứa ít hơn thùng $A$ là $9\text{ L}$. Thùng $B$ chứa bao nhiêu lít nước?
\choice
{$24\text{ L}$}
{$25\text{ L}$}
{$27\text{ L}$}
{\True $26\text{ L}$}
\loigiai{
$35 - 9 = 26\text{ (L)}$.
\textbf{Đáp án của câu hỏi là D}
}
\end{ex}

\begin{ex} %[2H3V18]
Một con lợn nặng $45\text{ kg}$, một con dê nhẹ hơn con lợn $17\text{ kg}$. Cả hai con nặng bao nhiêu ki-lô-gam?
\choice
{$71\text{ kg}$}
{\True $73\text{ kg}$}
{$75\text{ kg}$}
{$68\text{ kg}$}
\loigiai{
Con dê nặng: $45 - 17 = 28\text{ (kg)}$. Cả hai con nặng: $45 + 28 = 73\text{ (kg)}$.
\textbf{Đáp án của câu hỏi là B}
}
\end{ex}

\begin{ex} %[2H3V18]
Bác Lan có $30\text{ L}$ mật ong, bác rót vào các can loại $5\text{ L}$. Hỏi bác Lan rót được bao nhiêu can?
\choice
{\True $6$ can}
{$5$ can}
{$7$ can}
{$4$ can}
\loigiai{
$30 : 5 = 6$ (can).
\textbf{Đáp án của câu hỏi là A}
}
\end{ex}

\begin{ex} %[2H3V18]
Một bình thủy tinh nặng $500\text{ g}$, khi đựng đầy nước thì bình nặng $1\text{ kg } 500\text{ g}$. Lượng nước trong bình nặng bao nhiêu?
\choice
{$2\text{ kg}$}
{$500\text{ g}$}
{\True $1\text{ kg}$}
{$800\text{ g}$}
\loigiai{
Lượng nước nặng: $1\text{ kg } 500\text{ g} - 500\text{ g} = 1\text{ kg}$.
\textbf{Đáp án của câu hỏi là C}
}
\end{ex}

\begin{ex} %[2H3V18]
Tính giá trị: $15\text{ L} + 18\text{ L} - 9\text{ L}$.
\choice
{$22\text{ L}$}
{$26\text{ L}$}
{$25\text{ L}$}
{\True $24\text{ L}$}
\loigiai{
$15 + 18 = 33\text{ L}$; $33 - 9 = 24\text{ L}$.
\textbf{Đáp án của câu hỏi là D}
}
\end{ex}

\begin{ex} %[2H3V18]
Trong kho có $85\text{ kg}$ đỗ xanh. Ngày thứ nhất xuất $28\text{ kg}$, ngày thứ hai xuất $34\text{ kg}$. Trong kho còn lại bao nhiêu ki-lô-gam đỗ xanh?
\choice
{$25\text{ kg}$}
{\True $23\text{ kg}$}
{$21\text{ kg}$}
{$27\text{ kg}$}
\loigiai{
Tổng xuất: $28 + 34 = 62\text{ (kg)}$. Còn lại: $85 - 62 = 23\text{ (kg)}$.
\textbf{Đáp án của câu hỏi là B}
}
\end{ex}

\begin{ex} %[2H3V18]
Hai can chứa tất cả $40\text{ L}$ nước. Can thứ nhất chứa $18\text{ L}$. Can thứ hai chứa nhiều hơn can thứ nhất bao nhiêu lít nước?
\choice
{\True $4\text{ L}$}
{$5\text{ L}$}
{$3\text{ L}$}
{$6\text{ L}$}
\loigiai{
Can 2 chứa: $40 - 18 = 22\text{ (L)}$. Can 2 hơn can 1: $22 - 18 = 4\text{ (L)}$.
\textbf{Đáp án của câu hỏi là A}
}
\end{ex}

\begin{ex} %[2H3V18]
Một con ngỗng nặng $6\text{ kg}$, con vịt nhẹ hơn con ngỗng $3\text{ kg}$. Hỏi con ngỗng nặng gấp mấy lần con vịt?
\choice
{$3$ lần}
{$4$ lần}
{\True $2$ lần}
{$1$ lần}
\loigiai{
Con vịt nặng: $6 - 3 = 3\text{ (kg)}$. Con ngỗng nặng gấp con vịt: $6 : 3 = 2$ (lần).
\textbf{Đáp án của câu hỏi là C}
}
\end{ex}

\begin{ex} %[2H3V18]
Rót $16\text{ L}$ nước vào $4$ bình như nhau. Mỗi bình có bao nhiêu lít nước?
\choice
{$3\text{ L}$}
{$5\text{ L}$}
{$6\text{ L}$}
{\True $4\text{ L}$}
\loigiai{
$16 : 4 = 4\text{ (L)}$.
\textbf{Đáp án của câu hỏi là D}
}
\end{ex}

\textbf{Phần 2. Trắc nghiệm đúng sai}
\setcounterpageref{ex}{0}

\begin{ex} %[2H3V18]
Xét tính đúng sai của các bài toán tổng hợp về đo lường:
\choiceTF
{\True Tổng của 25 kg và 35 kg bằng 60 kg.}
{Phép tính 45 L - 18 L có kết quả là 37 L.}
{\True Một can 10 L đựng được nhiều nước hơn hai can 4 L gộp lại.}
{Khối lượng luôn đo bằng Lít còn dung tích luôn đo bằng Ki-lô-gam.}
\loigiai{
\begin{itemize}
    \item Ý a) Đúng: $25 + 35 = 60\text{ kg}$.
    \item Ý b) Sai: $45 - 18 = 27\text{ L}$ (không phải 37 L).
    \item Ý c) Đúng: Hai can 4 L là $4 + 4 = 8\text{ L} < 10\text{ L}$. Đúng.
    \item Ý d) Sai: Ngược lại: khối lượng đo bằng kg, dung tích đo bằng Lít.
\end{itemize}
\begin{itemize}
    \item a) \textbf{Đúng}.
    \item b) \textbf{Sai}.
    \item c) \textbf{Đúng}.
    \item d) \textbf{Sai}.
\end{itemize}
}
\end{ex}

\begin{ex} %[2H3V18]
Một cửa hàng có 3 thùng sữa: Thùng 1 có 24 L, Thùng 2 có 28 L, Thùng 3 có 32 L:
\choiceTF
{\True Cả ba thùng có tất cả 84 L sữa.}
{\True Thùng 3 nhiều hơn Thùng 1 là 8 L sữa.}
{Thùng 2 ít hơn Thùng 3 là 6 L sữa.}
{\True Trung bình số sữa của ba thùng là 28 L.}
\loigiai{
\begin{itemize}
    \item Ý a) Đúng: $24 + 28 + 32 = 84\text{ L}$.
    \item Ý b) Đúng: $32 - 24 = 8\text{ L}$.
    \item Ý c) Sai: $32 - 28 = 4\text{ L}$ (không phải 6 L).
    \item Ý d) Đúng: $(24 + 28 + 32) : 3 = 84 : 3 = 28\text{ L}$. Đúng.
\end{itemize}
\begin{itemize}
    \item a) \textbf{Đúng}.
    \item b) \textbf{Đúng}.
    \item c) \textbf{Sai}.
    \item d) \textbf{Đúng}.
\end{itemize}
}
\end{ex}

\textbf{Phần 3. Trả lời ngắn}
\setcounterpageref{ex}{0}

\begin{ex} %[2H3V18]
Một người mua 18 kg gạo tẻ và 9 kg gạo nếp. Hỏi người đó mua tất cả bao nhiêu ki-lô-gam gạo?
\shortans{27}
\loigiai{
\textbf{Đáp số:} $27$\\
$18 + 9 = 27\text{ kg}$.
}
\end{ex}

\begin{ex} %[2H3V18]
Bể thứ nhất chứa 45 lít nước, bể thứ hai chứa 38 lít nước. Hỏi bể thứ nhất chứa nhiều hơn bể thứ hai bao nhiêu lít nước?
\shortans{7}
\loigiai{
\textbf{Đáp số:} $7$\\
$45 - 38 = 7\text{ L}$.
}
\end{ex}

\begin{ex} %[2H3V18]
Tính giá trị: $54\text{ kg} - 28\text{ kg} + 14\text{ kg}$.
\shortans{40}
\loigiai{
\textbf{Đáp số:} $40$\\
$54 - 28 = 26$; $26 + 14 = 40\text{ kg}$.
}
\end{ex}

\begin{ex} %[2H3V18]
Một can dầu chứa đầy được 15 lít. Người ta rót ra 8 lít thì trong can còn lại bao nhiêu lít dầu?
\shortans{7}
\loigiai{
\textbf{Đáp số:} $7$\\
$15 - 8 = 7\text{ L}$.
}
\end{ex}

\begin{ex} %[2H3V18]
Mẹ nuôi một đàn gà ăn hết 25 kg thóc trong tháng đầu và 28 kg thóc trong tháng thứ hai. Hỏi cả hai tháng đàn gà ăn hết bao nhiêu ki-lô-gam thóc?
\shortans{53}
\loigiai{
\textbf{Đáp số:} $53$\\
$25 + 28 = 53\text{ kg}$.
}
\end{ex}

\begin{ex} %[2H3V18]
Có 32 lít nước mắm chia đều vào các can 4 lít. Hỏi chia được bao nhiêu can?
\shortans{8}
\loigiai{
\textbf{Đáp số:} $8$\\
$32 : 4 = 8$ can.
}
\end{ex}
